\begin{abstract}
We propose two search systems for the University of Calgary course timetabling problem: assigning lectures and tutorials to weekly time slots while satisfying hard constraints and minimizing soft-constraint penalties. Both systems begin with OR-tree search, whose nodes describe partial timetables. Depth-first exploration with minimum remaining values selects offerings and explores legal slot assignments until a complete valid timetable is found or all alternatives fail.

The systems differ in how they improve this initial timetable. Proposal A uses set-based search, representing a timetable as a set of assignment facts. Greedy hill climbing selects Move or Swap rules that preserve hard constraint validity while reducing the overall penalty value. Proposal B uses AND-tree search to divide the optimization problem into assignment cases, applying branch-and-bound to prune cases that cannot improve the best timetable found.

Both systems retain a valid timetable during optimization. Proposal A may stop at a local optimum, whereas Proposal B can establish global optimality when its search completes. We define each model, its search control and stopping conditions, and demonstrate rule applications through worked examples.
\end{abstract}